%==========================================================================================
\section{Follow-up (22 October): where the remaining error comes from, and five attempts to reduce it}\label{sec:round4}

This section asks why the fair-horizon model (\S\ref{sec:fairh8}) still trails the U-Net,
and reports five attempts to fix it. All attempts keep the requirements: exactly linear
latent dynamics, continuous time, Mach conditioning, forecasting-only end-to-end training,
and the 2 $\to$ 8 horizon. Details are in \texttt{EXPERIMENT\_RESULTS.md} \S14 and
\texttt{LITERATURE\_REVIEW\_V3.md}.

\subsection{Diagnosis}\label{sec:diag4}

\textbf{The first-step error is all reconstruction.} I split the one-step error on the test
set into two parts: encoding and decoding the current frame, and the error added by one
linear time step (\texttt{scripts/lead1\_split.py}).

\begin{center}\small
\begin{tabular}{lccc}
\toprule
One-step error, interp ($\times10^{-3}$) & reconstruction & added by the linear step & total\\
\midrule
fair 2 $\to$ 8, latent 256 & 0.317 & 0.008 & 0.298\\
fair 2 $\to$ 8, latent 128 & 0.382 & 0.010 & 0.361\\
ACDM-setting training (2 $\to$ 1) & 0.092 & 0.022 & 0.088\\
\bottomrule
\end{tabular}
\end{center}

\textbf{But the first step explains only part of the gap.} Gap in MSE to each baseline, by
lead (interp):

\begin{center}\small
\begin{tabular}{lcccccc}
\toprule
 & 60-step average & lead 1 & lead 10 & lead 20 & lead 30 & lead 58\\
\midrule
vs ACDM  & 0.47 & 0.23 & 0.18 & 0.54 & 0.92 & 0.03\\
vs U-Net & 1.96 & 0.24 & 0.35 & 1.13 & 2.16 & 3.58\\
\bottomrule
\end{tabular}
\end{center}

\begin{itemize}
\item Against ACDM, the reconstruction error is about half of the gap.
\item Against the U-Net it is about one eighth. The rest grows with lead time.
\end{itemize}

\textbf{What grows.} I projected the error onto the true trajectory's own POD modes
(\texttt{scripts/phase\_drift.py}).
\begin{itemize}
\item The shedding frequency is accurate: within 0.2--1.5\% on extrapolation, against
  ACDM 2.7\% and U-Net 1.2\%.
\item Excess error over the U-Net at lead 30:
  \begin{itemize}
  \item 49\% lies outside the span of the true trajectory (blurred or misplaced structure);
  \item 32\% is in POD modes 3--10 (harmonics and secondary modes);
  \item 8\% is in the shedding pair.
  \end{itemize}
\item At lead 58 the shedding pair carries 33\% (phase error 26--37$^\circ$ vs 14$^\circ$ for
  the U-Net), and the outside-span part 53\%.
\end{itemize}

\textbf{Checks on the trained model, at inference only} (\texttt{scripts/infer\_ablate.py}).
\begin{itemize}
\item \textbf{Damping helps.} Scaling every oscillator's damping by 1 / 0.5 / 0.25 / 0
  gives interp 3.43 / 3.58 / 3.68 / 3.81. The damping hedges uncertain components, and the
  modes that carry most of the latent are already almost neutral.
\item \textbf{The dynamics contain no harmonics.} At most 0.5\% of the latent's temporal
  variance sits at or above twice the shedding frequency. Harmonics of the wake come only
  from the decoder.
\end{itemize}

\textbf{The flow is only partly periodic} (training data only;
\texttt{scripts/shedding\_freq.py}).
\begin{itemize}
\item Below Ma 0.875 the shedding frequency is $2.45\pm0.16$ per time unit. It scatters by
  6.5\% between neighbouring Mach numbers.
\item Ma 0.88--0.89 is a separate regime at about 6.3.
\item On average only 57\% of the dominant mode's spectrum lies on the harmonic lattice.
\end{itemize}

\subsection{Literature}\label{sec:lit4}

Two source-verified reviews covered about 47 papers.
\begin{itemize}
\item \textbf{Most remedies target a problem this model does not have.} Pushforward, input
  noise, unrolling and test-time correction all target autoregressive drift. Here every lead
  is one exact matrix exponential from true input frames.
\item \textbf{Evidence on Tra itself} (our re-scoring of the official checkpoints):
  \begin{itemize}
  \item spectral refinement (PDE-Refiner) scores 7.50 / 5.61;
  \item every latent model in the ACDM paper scores 6.4--7.4 on interpolation, so the
    CT-KAE is already the best latent model on this benchmark;
  \item only the stabilised pixel-space U-Nets are ahead (1.53 / 1.66 unrolled, 1.83 / 1.42
    with input noise).
  \end{itemize}
\item \textbf{The two constraint-compatible ideas both reviews ranked highest:}
  \begin{itemize}
  \item a harmonic-lattice generator (Bagheri 2013; Lange et al.\ 2021);
  \item multi-start temporal consistency (tcKAE, Nayak et al.\ 2024).
  \end{itemize}
  No paper reports a linear-latent model reaching the stabilised U-Net under comparable
  constraints.
\end{itemize}

\subsection{Five attempts}\label{sec:attempts4}

All runs used seed 0 and \texttt{last.pth}. The success criterion was fixed before the
runs: interp $\le 3.19$ (two seed standard deviations below the plain model's mean) with
extrap $\le 2.29$. DA is the range over the six df = 9 sweeps.

\begin{center}\footnotesize
\begin{tabular}{lcccccc}
\toprule
Seed 0 & interp & extrap & long & lead 1 & lead 58 & DA df = 9\\
\midrule
\textbf{fair 2 $\to$ 8, latent 256 (reference)} & \textbf{3.43} & \textbf{1.67} & 16.80 & 0.261 & \textbf{8.17} & \textbf{0.091--0.095}\\
+ detail latent, advection        & 8.85 & 5.53 & 16.98 & \textbf{0.021} & 15.8 & 0.142--0.176\\
+ detail latent, decay only       & 16.44 & 9.80 & 24.52 & 0.026 & 18.3 & 0.199--0.217\\
+ detail latent, advection, decay floor & 12.79 & 5.99 & 24.73 & 0.030 & 21.5 & 0.292--0.306\\
harmonic lattice (64 of 128 pairs) & 3.87 & 2.29 & \textbf{15.35} & 0.272 & 9.77 & 0.154--0.163\\
multi-start consistency           & 4.79 & 2.71 & 16.32 & 0.217 & 12.1 & 0.099--0.104\\
\bottomrule
\end{tabular}
\end{center}

\noindent\textbf{Result: none met the criterion.}
\begin{itemize}
\item \textbf{Detail latent.}
  \begin{itemize}
  \item \emph{Design:} a spatial latent with a Fourier-diagonal advection and decay generator.
    It is exactly linear, which is tested against the dense matrix exponential to $10^{-9}$.
  \item \emph{What it fixes:} the first step, to 0.02--0.03, as good as ACDM and the U-Net.
  \item \emph{What goes wrong:} the model moves the near-term dynamics into the detail
    latent. The oscillators go static or leave the shedding frequency, and the decoder never
    learns to forecast without the detail. Beyond the 8-frame window the forecast collapses,
    and a decay floor does not prevent it.
  \end{itemize}
\item \textbf{Harmonic lattice.}
  \begin{itemize}
  \item The long rollout improves slightly; 15.35 is below all three reference seeds.
  \item Interpolation and DA get worse.
  \item This is consistent with the irregular, partly non-periodic shedding found above.
  \end{itemize}
\item \textbf{Multi-start consistency.}
  \begin{itemize}
  \item The first step improves (0.217), and so do the DA tuning error and the encode-init
    floor.
  \item The 60-step forecast and the df = 9 sweeps get worse.
  \end{itemize}
\end{itemize}

\subsection{Conclusion}\label{sec:concl4}

\begin{itemize}
\item \textbf{The binding constraint is the 8-frame supervised window.} Every change that
  makes the near term sharper also gives the model a way to do well within 8 frames that
  does not carry beyond them. Nothing in a forecasting-only, 8-frame objective penalises
  that.
\item \textbf{The plain fair-horizon model remains the best} of everything tried. It is the
  best latent model on Tra, level with ACDM on extrapolation, ahead of it on the 240-step
  rollout, and about 2.7$\times$ better than ACDM and the U-Net at df = 9 data
  assimilation.
\item \textbf{What would still move the gap to the U-Net.} The one remaining lever with a
  clear mechanism is relaxing the horizon rule: the same 8 frames spread over a longer
  physical time span, which constrains frequencies and phases better. That is a decision
  about the comparison protocol, not a modelling fix.
\end{itemize}
